\title{DeepSeekMath: Pushing the Limits of Mathematical Reasoning in Open Language Models}

\begin{abstract}
Mathematical reasoning remains a formidable obstacle for language models due to its intrinsic complexity and structured demands. This study presents DeepSeekMath~7B, developed by extending the pre-training of DeepSeek-Coder-Base-v1.5 7B utilizing 120B mathematics-focused tokens extracted from Common Crawl, as well as data in natural language and code formats. Remarkably, DeepSeekMath~7B attains a 51.7\% score on the highly challenging MATH benchmark without dependence on external toolkits or voting schemes, reaching the competence level seen in Gemini-Ultra and GPT-4. When employing self-consistency across 64 generated samples, DeepSeekMath~7B further elevates its performance to 60.9\% on the same benchmark. The notable improvement in mathematical reasoning for DeepSeekMath~is underpinned by two main factors: first, the effective exploitation of public web data enabled by a carefully crafted data filtering and selection process; second, the development of Group Relative Policy Optimization (GRPO)—an adaptation of Proximal Policy Optimization (PPO)—which bolsters both reasoning proficiency and memory efficiency during optimization.
\end{abstract}

\section{Introduction}

Large language models (LLM) have fundamentally transformed AI approaches to mathematical reasoning, catalyzing progress on benchmarks evaluating both quantitative \citep{MATH} and geometric reasoning abilities \citep{trinh2024solving}. These models also provide substantial support to humans addressing intricate mathematical questions \citep{tao}. Nevertheless, leading-edge solutions such as GPT-4 \citep{gpt4} and Gemini-Ultra \citep{gemini} are not openly accessible, and the performance of existing open-source alternatives falls significantly short in comparison.

This paper presents DeepSeekMath, a language model specialized in mathematics that demonstrates superior performance compared to existing open-source models and closely rivals GPT-4 in academic evaluations. To support the development of this model, we construct the DeepSeekMath Corpus, a substantial and high-quality dataset encompassing 120B math tokens. The data collection process relies on extracting relevant information from the Common Crawl (CC) using a classifier built on the fastText framework \citep{joulin2016fasttext}. For its initial training, this classifier uses examples from OpenWebMath \citep{openwebmath} as positive cases and samples from various unrelated web pages as negative cases. The trained classifier is subsequently applied to the CC to identify further positive samples, which undergo additional filtering and verification via human annotation. These refined samples are then incorporated to retrain and enhance the classifier. Our empirical analysis demonstrates that the resulting dataset maintains high quality: the DeepSeekMath-Base 7B model attains a score of 64.2\% on GSM8K \citep{gsm8k} and 36.2\% on the MATH dataset \citep{MATH}, surpassing the performance of Minerva 540B \citep{minerva}. The DeepSeekMath Corpus is inherently multilingual, resulting in observed gains on Chinese mathematical tasks \citep{wei2023cmath,agieval}. We propose that our methodology for curating mathematical datasets offers a valuable foundation for future research, while recognizing that there remains considerable potential for further advancement.

DeepSeekMath-Base is built upon DeepSeek-Coder-Base-v1.5 7B \citep{deepseek-coder}, as our findings indicate that initializing with a code-centric model yields better outcomes than using a generic LLM. Additionally, we find that training the model on mathematical data not only strengthens its mathematical proficiency but also leads to notable improvements on general benchmarks like MMLU \citep{mmlu} and BBH \citep{bbh}, thereby enhancing overall reasoning skills.

Following the pre-training stage, we fine-tune DeepSeekMath-Base on mathematical instruction datasets, leveraging approaches such as chain-of-thought \citep{cot}, program-of-thought \citep{pot,pal}, and reasoning with tool integration \citep{tora}. This instruction-tuned variant, DeepSeekMath-Instruct 7B, outperforms all other 7B-sized models and reaches parity with leading instruction-tuned open-source models of 70B parameters.

Moreover, we propose Group Relative Policy Optimization (GRPO), a novel reinforcement learning algorithm derived from Proximal Policy Optimization (PPO) \citep{schulman2017proximal}. GRPO eliminates the need for a critic by instead using group-based score baselines, resulting in marked efficiency gains during training. Utilizing only a portion of English-language instruction-tuning data, GRPO brings substantial gains over the robust DeepSeekMath-Instruct, achieving significant improvements on both in-domain (GSM8K: 82.9\% $\rightarrow$ 88.2\%, MATH: 46.8\% $\rightarrow$ 51.7\%) and out-of-domain mathematical evaluations (e.g., CMATH: 84.6\% $\rightarrow$ 88.8\%) during RL-based training. We also present a unified perspective encompassing methods such as Rejection Sampling Fine-Tuning (RFT) \citep{yuan2023scaling}, Direct Preference Optimization (DPO) \citep{dpo}, PPO, and GRPO, framing them as manifestations of either direct or streamlined RL methodologies. Extensive empirical analyses, including comparisons between online and offline regimes, outcome-oriented versus process supervision, and single-step versus iterative RL, are conducted to unpack the crucial aspects of this framework. Finally, we elucidate the mechanisms by which our RL strategies augment instruction-tuned models and highlight future avenues for enhancing RL methods within this consolidated paradigm.

\subsection{Contributions}
We present advancements in large-scale mathematical pre-training, together with an in-depth study of reinforcement learning techniques.

\textbf{Math Pre-Training at Scale}

\begin{itemize}[topsep=0pt]
    \item Through our investigation, we demonstrate that the widely available Common Crawl dataset holds substantial resources for mathematical data. Utilizing a carefully constructed data selection pipeline, we develop the DeepSeekMath Corpus—a refined collection of 120B tokens sourced from mathematically relevant web pages. This corpus surpasses the math-related datasets used by Minerva \citep{minerva} by nearly sevenfold, and is approximately nine times larger than the newly introduced OpenWebMath \citep{openwebmath}.

    \item The DeepSeekMath-Base 7B, our pre-trained foundational model, reaches a level of mathematical performance similar to Minerva 540B \citep{minerva}. This outcome suggests that sheer model size is not the exclusive factor underlying success in mathematical reasoning; instead, robust results can be obtained from smaller models if pre-trained on superior-quality data.

    \item We report insights gained from various math-focused training regimens. Specifically, training models on code before mathematical tasks significantly enhances their ability to address mathematical problems, regardless of whether external tools are utilized. This partially addresses the ongoing debate: \textit{does exposure to code training improve reasoning skills?} Our evidence indicates a positive impact, especially in the domain of mathematical reasoning.

    \item Contrary to prevalent practices in the field, including arXiv papers during pre-training yields no discernible gains across all mathematical evaluation benchmarks considered in this work.
\end{itemize}

\textbf{Exploration and Analysis of Reinforcement Learning}

\begin{itemize}[topsep=0pt]
    \item We present Group Relative Policy Optimization (GRPO), a reinforcement learning algorithm designed for both efficiency and performance. GRPO eliminates the need for a critic network by leveraging group-based score baselines, thus substantially decreasing the computational requirements relative to Proximal Policy Optimization (PPO).
    \item We show that applying GRPO markedly improves the capabilities of our instruction-tuned model DeepSeekMath-Instruct, using only the instruction-tuning corpus. Additionally, we observe that GRPO facilitates generalization, as evidenced by performance gains on out-of-domain evaluation sets during the reinforcement learning stage.
    \item We articulate a comprehensive framework that integrates a variety of approaches, including RFT, DPO, PPO, and GRPO. Through a range of empirical studies—such as comparisons between online and offline training, outcome versus process-based supervision, and analyses of single-turn versus iterative reinforcement learning—we examine the foundational components and underlying mechanisms of this paradigm.
    \item Guided by this framework, we investigate the sources of reinforcement learning effectiveness and highlight several promising avenues to further enhance reinforcement learning strategies for large language models (LLMs).
\end{itemize}

\subsection{Summary of Evaluations and Metrics}

\begin{itemize}[topsep=0pt]
    \item \textbf{English and Chinese Mathematical Reasoning}: Our evaluation protocol encompasses a wide array of English and Chinese benchmarks, testing mathematical reasoning from elementary through tertiary education. English evaluations are performed on GSM8K \citep{gsm8k}, MATH \citep{MATH}, SAT \citep{llemma}, OCW Courses \citep{minerva}, and MMLU-STEM \citep{mmlu}, while the Chinese tasks involve MGSM-zh \citep{mgsm}, CMATH \citep{wei2023cmath}, Gaokao-MathCloze \citep{agieval}, and Gaokao-MathQA \citep{agieval}. Our assessment covers the generation of full, self-contained textual solutions (without access to tools), as well as code-based problem solving in Python.

    On English datasets, DeepSeekMath-Base shows competitive performance relative to Minerva 540B \citep{minerva}, which is a closed-source model, and it consistently outperforms all available open-source baselines—such as Mistral 7B \citep{mistral} and Llemma-34B \citep{llemma}—regardless of math pre-training status, often by a notable margin. Remarkably, on Chinese datasets, DeepSeekMath-Base outperforms previous baselines, plausibly due to our decision not to limit math pre-training data to English sources as in prior work \citep{minerva,llemma}, but to deliberately include high-quality data in other languages as well. Leveraging instruction-tuning and reinforcement learning, both DeepSeekMath-Instruct and DeepSeekMath-RL deliver robust results, surpassing 50\% accuracy on the highly challenging MATH competition dataset—a milestone achieved for the first time in open-source models.
\end{itemize}

\item \textbf{Formal Mathematics}: The informal-to-formal theorem proving benchmark introduced in \citep{dsp_proof} and based on miniF2F \citep{minif2f}, with Isabelle \citep{isabelle} serving as the designated proof assistant, is used to assess DeepSeekMath-Base. On this task, DeepSeekMath-Base exhibits notable performance in few-shot autoformalization scenarios.
\item \textbf{Natural Language Understanding, Reasoning, and Code}: To thoroughly characterize the general capabilities of models in understanding, reasoning, and coding, we assess DeepSeekMath-Base using the Massive Multitask Language Understanding (MMLU) benchmark \citep{mmlu}, which spans 57 multiple-choice subject areas, BIG-Bench Hard (BBH) \citep{bbh}, a suite of 23 primarily multi-step reasoning tasks, as well as HumanEval \citep{codex} and MBPP \citep{mbpp}, both standard evaluations for code-generating language models. Pre-training on mathematical data leads to measurable gains in both language comprehension and reasoning.

\section{Math Pre-Training}

\subsection{Data Collection and Decontamination} This section describes the methodology for assembling the DeepSeekMath Corpus from the Common Crawl dataset. As shown in Figure \ref{fig:data_collect}, an iterative pipeline is utilized, beginning with a curated, high-quality set of math-related data (the seed corpus) and systematically expanding it to harvest a large-scale mathematical dataset from Common Crawl. Although designed for mathematics, this method can extend to other specialized fields, such as programming.

We start the process with OpenWebMath \citep{openwebmath}, a repository of well-vetted mathematical web documents, serving as the seed corpus. A fastText model \citep{joulin2016fasttext} is trained to identify additional web pages similar to those found in OpenWebMath. For training purposes, 500,000 examples from the seed corpus are labeled as positive samples, while another 500,000 web pages from Common Crawl are sampled as negatives. We use an open-source fastText implementation\footnote{\url{https://fasttext.cc}}, setting the vector size to 256, learning rate to 0.1, word n-gram maximum length to 3, minimum frequency of word occurrence to 3, and the number of training epochs to 3. In order to downsize the vast Common Crawl dataset, both URL-based and near-deduplication strategies are applied, reducing it to 40 billion HTML pages. Mathematical content is then retrieved from this deduplicated set using the fastText model. To further ensure the quality of the extracted mathematical content, the pages are ranked by their fastText model prediction scores, and only those with the highest rankings are retained. We empirically evaluate the size of retained data by running pre-training experiments with the top 40B, 80B, 120B, and 160B tokens; for the initial iteration, we select the top 40B tokens.

Following the initial phase of data collection, a significant number of mathematical web pages remain uncaptured. This is primarily attributed to the limited diversity within the positive examples used to train the fastText model. To address this limitation and enhance the model's effectiveness, we augment the seed corpus with additional mathematical web resources. Specifically, we partition the complete Common Crawl dataset into separate domains, with each domain comprising web pages that share an identical base URL. For every domain, we determine the proportion of web pages successfully collected during the first round. Domains in which more than 10\% of the web pages have been gathered are classified as math-focused (for instance, \url{mathoverflow.net}). 

Within these math-related domains, we perform manual annotation to identify URLs that correspond to mathematical content (e.g., \url{mathoverflow.net/questions}). Uncollected web pages connected to these URLs are subsequently incorporated into the seed corpus, yielding a richer set of positive samples. This procedure strengthens the fastText model, allowing it to retrieve a broader range of mathematical web content in later collection cycles. After completing four rounds of data collection, we accumulate 35.5M mathematical web pages, amassing a total of 120B tokens. In the fourth round, we observe that approximately 98\% of the material was already obtained in the preceding iteration, prompting us to halt further collection.

In order to prevent contamination from benchmark datasets, we apply a filtering process as outlined by \cite{deepseek-coder}. Web pages containing text from well-known English mathematical benchmarks, such as GSM8K \citep{gsm8k} and MATH \citep{MATH}, as well as Chinese datasets like CMATH \citep{wei2023cmath} and AGIEval \citep{agieval}, are excluded. The filtering mechanism removes from our math training corpus any text segment featuring a 10-gram sequence that matches exactly with a substring from the evaluation benchmarks. For cases in which the benchmark text is under 10 grams but at least 3 grams, precise matching is employed to discard any contaminated web pages.

\subsection{Validating the Quality of the DeepSeekMath~Corpus}

We conduct pre-training experiments to evaluate the DeepSeekMath~Corpus relative to recently developed mathematical text corpora:

\begin{itemize}[topsep=0pt]
    \item \textbf{MathPile} \citep{mathpile}: A multi-source dataset (8.9B tokens), compiled from sources including textbooks, Wikipedia, ProofWiki, CommonCrawl, StackExchange, and arXiv—with more than 85\% of the data derived from arXiv;
    \item \textbf{OpenWebMath} \citep{openwebmath}: Data from CommonCrawl, filtered for mathematical relevance, containing a total of 13.6B tokens;
    \item \textbf{Proof-Pile-2} \citep{llemma}: A mathematical corpus incorporating OpenWebMath, AlgebraicStack (10.3B tokens of mathematical code), and arXiv articles (28.0B tokens). In our experiments on Proof-Pile-2, we adhere to the arXiv:Web:Code ratio of 2:4:1 prescribed by \cite{llemma}.
\end{itemize}

\subsubsection{Training Setting}
\label{sec:quality-policy}
We perform mathematical specialization on a foundational language model comprising 1.3 billion parameters, utilizing the same architecture as DeepSeek LLMs \citep{deepseek-llm}, referred to as DeepSeek-LLM 1.3B. Each mathematical dataset serves as the basis for a distinct model, trained for a total of 150 billion tokens. All training runs leverage the HAI-LLM \citep{haillm} framework, which is noted for its efficiency and minimal overhead. Consistent with established practices in DeepSeek LLM training, we adopt the AdamW optimizer \citep{adamW} configured with $\beta_1=0.9$, $\beta_2=0.95$, and a $\mathrm{weight\_decay}$ of 0.1. The learning rate schedule follows a multi-step protocol: it rises to its maximum after 2,000 warmup steps, is reduced to 31.6\% of its peak by the time 80\% of total training has been completed, and further declines to 10.0\% by the 90\% mark. The highest learning rate is fixed at 5.3e-4. We employ a batch size of 4 million tokens and a maximum sequence length of 4,000 tokens.

\subsubsection{Evaluation Results}

\textbf{The DeepSeekMath~Corpus stands out for its high quality, multilingual breadth, and unmatched scale among mathematical corpora.}

\begin{itemize}[topsep=0pt]
    \item \textbf{High-quality}:
    To assess downstream capabilities, we evaluate performance on 8 mathematical tasks employing few-shot chain-of-thought prompting as described in \cite{cot}. According to Table \ref{tab:corpora_comparison}, models trained on the DeepSeekMath~Corpus exhibit clear superiority over alternatives. Moreover, Figure \ref{fig:corpora_comparisons} demonstrates that, after training on 50B tokens (one complete pass over Proof-Pile-2), the model utilizing DeepSeekMath outperforms its Proof-Pile-2 counterpart, underscoring the superior overall quality of the DeepSeekMath dataset.
    \item \textbf{Multilingual}:
    The DeepSeekMath~Corpus includes material in a variety of languages, with English and Chinese as the principal representatives. Table \ref{tab:corpora_comparison} demonstrates that this diversity benefits the model’s mathematical reasoning abilities in both English and Chinese. By contrast, preexisting corpora that focus heavily on English not only provide limited gains but can sometimes negatively impact Chinese-language mathematical tasks.
    \item \textbf{Large-scale}:
    DeepSeekMath~Corpus far exceeds previous mathematical datasets in terms of scale. As illustrated in Figure \ref{fig:corpora_comparisons}, the DeepSeek-LLM 1.3B trained on DeepSeekMath experiences more rapid progress and sustained improvements over time. Competing corpora, due to their smaller size, are quickly cycled through during training, leading models trained on them to plateau more rapidly.
\end{itemize}

\subsection{Training and Evaluating DeepSeekMath-Base 7B}

This section details the development of DeepSeekMath-Base 7B, a foundational model designed for advanced mathematical reasoning. The initial weights are sourced from DeepSeek-Coder-Base-v1.5 7B \citep{deepseek-coder}, followed by training on a total of 500B tokens. The data utilized for this process comprises 56\% DeepSeekMath Corpus, 4\% AlgebraicStack, 10\% arXiv, 20\% code from Github, and 10\% natural language text drawn from Common Crawl, with content in both English and Chinese. The majority of hyperparameters align with those outlined in Section \ref{sec:quality-policy}; however, the maximum learning rate is set to 4.2e-4, and training uses a batch size of 10M tokens.

To rigorously benchmark the mathematical proficiency of DeepSeekMath-Base 7B, we assess its capacity to independently generate comprehensive mathematical solutions, leverage external tools when necessary, and engage in formal theorem proving. Additionally, we examine broader aspects of model competence, encompassing abilities in natural language comprehension, logical reasoning, and software development tasks.

\paragraph{Mathematical Problem Solving with Step-by-Step Reasoning}

The model's mathematical problem-solving skills are tested via few-shot chain-of-thought prompting \citep{cot}, spanning eight diverse English and Chinese evaluation datasets. These include benchmarks on quantitative reasoning such as GSM8K \citep{gsm8k}, MATH \citep{MATH}, and CMATH \citep{wei2023cmath}, as well as assessments focused on multiple-choice formats like MMLU-STEM \citep{mmlu} and Gaokao-MathQA \citep{agieval}. Collectively, these cover mathematical domains ranging from elementary through collegiate difficulty.

As detailed in Table \ref{tab:base_math_cot}, DeepSeekMath-Base 7B achieves top-tier results on all eight benchmarks compared to other open-source base models, which include the popular general-purpose Mistral 7B \citep{mistral} and the more recent Llemma 34B \citep{llemma}, both trained on specialized mathematical corpora such as Proof-Pile-2 \citep{llemma}. Particularly impressive is its performance on the MATH dataset, designed for mathematical competitions, where DeepSeekMath-Base outperforms all available open-source base models by an absolute margin of more than 10\%. Furthermore, it even exceeds the results of Minerva 540B \citep{minerva}—a much larger, proprietary model (77 times the parameter size), which itself is based on PaLM \citep{palm} with supplemental mathematical training.

\paragraph{Mathematical Problem Solving with Tool Use} For tasks involving program-supported mathematical reasoning, we conduct evaluations on GSM8K and MATH datasets using a few-shot program-of-thought prompting framework \citep{pot,pal}. Models are required to address each question by composing a Python script, with access to modules such as \textit{math} and \textit{sympy} to handle complex calculations. The program’s output serves as the model's answer. Table \ref{tab:base_math_pot_proof} reveals that DeepSeekMath-Base 7B surpasses Llemma 34B, the previous best-performing model in this area.

\paragraph{Formal Mathematics}

Automation of formal proofs is increasingly valued for its potential to verify mathematical arguments and improve workflow efficiency. For assessment, we task DeepSeekMath-Base 7B with informal-to-formal proof generation as described in \citep{dsp_proof}—converting an informal problem statement, its formalized version, and a natural language proof into a rigorous formal proof. Our experiments utilize miniF2F \citep{minif2f}, a formalized mathematics challenge set at Olympiad difficulty, where the model is prompted with few examples to produce formal proofs in Isabelle. Following the approach from \cite{dsp_proof}, the model first generates an outline (proof sketch), after which the standard Sledgehammer prover \citep{sledgehammer} is used to automatically complete the proof details. The results in Table \ref{tab:base_math_pot_proof} highlight DeepSeekMath-Base 7B’s robust capabilities in automatic proof formalization.

\paragraph{Natural Language Understanding, Reasoning, and Code}

To assess capabilities in broader domains, we test for natural language understanding with MMLU \citep{mmlu}, logical reasoning with BBH \citep{bbh}, and programming proficiency using HumanEval \citep{codex} and MBPP \citep{mbpp}. According to Table \ref{tab:understanding_reasoning_code}, DeepSeekMath-Base 7B delivers substantial improvements over its predecessor, DeepSeek-Coder-Base-v1.5 \citep{deepseek-coder}, particularly on MMLU and BBH, evidencing the beneficial impact of mathematical pretraining on comprehension and reasoning tasks. Through continued exposure to code tokens, the model also preserves the coding abilities of DeepSeek-Coder-Base-v1.5 on the programming benchmarks. Across these three reasoning and coding evaluations, DeepSeekMath-Base 7B consistently outperforms the generalist Mistral 7B \citep{mistral}.

\section{Supervised Fine-Tuning}

\subsection{SFT Data Curation}

We develop a diverse instruction-tuning dataset for mathematics that includes both English and Chinese problems, spanning a wide range of mathematical disciplines and difficulty levels. Each problem is accompanied by a corresponding solution expressed using chain-of-thought (CoT) \citep{cot}, program-of-thought (PoT) \citep{pot,pal}, or tool-integrated reasoning formats \citep{tora}. The dataset contains a total of 776K training samples.

\begin{itemize}[topsep=0pt]
    \item \textbf{English mathematical datasets}: For the English portion, we enhance GSM8K and MATH problems by providing tool-integrated solution annotations. We further incorporate selected samples from MathInstruct \citep{MathInstruct} and include the training split from Lila-OOD \citep{lila}, where both CoT and PoT methods are used in solutions. This portion of the dataset spans numerous areas such as algebra, probability, number theory, calculus, and geometry.
    \item \textbf{Chinese mathematical datasets}: We gather K-12 mathematics questions in Chinese, which cover 76 distinct sub-fields including, for instance, linear equations. Each item features solutions prepared in both CoT and tool-integrated reasoning formats.
\end{itemize}

\subsection{Training and Evaluating DeepSeekMath-Instruct 7B}

This section presents DeepSeekMath-Instruct 7B, a model fine-tuned on mathematical instructions using DeepSeekMath-Base as a foundation. For training, we concatenate samples at random up to a 4K token context window. The model is trained for 500 iterations, employing a batch size of 256 and a fixed learning rate of 5e-5.

To assess mathematical capabilities, we evaluate the models—with and without tool usage—across four benchmarks focused on quantitative reasoning in both English and Chinese. Performance is compared with the foremost models available, including:

\begin{itemize}[topsep=0pt]
    \item \textbf{Closed-source models}: Our comparisons encompass (1) the GPT series, most notably GPT-4 \citep{gpt4} and the GPT-4 Code Interpreter~\footnote{\url{https://openai.com/blog/chatgpt-plugins\#\#code-interpreter}}, (2) Gemini Ultra and Pro \citep{gemini}, (3) Inflection-2 \citep{inflection-2}, (4) Grok-1~\footnote{\url{https://x.ai/model-card}}, as well as recently introduced models from major Chinese firms such as (5) Baichuan-3~\footnote{\url{https://www.baichuan-ai.com}} and (6) the new GLM-4~\footnote{\url{https://open.bigmodel.cn/dev/api\#glm-4}} from the GLM family \citep{glm}.
\end{itemize}

These models are designed for a wide range of applications, with the majority having been subjected to multiple rounds of alignment refinement.

\item \textbf{Open-source models} encompass both general-purpose and mathematically-enhanced systems. General-purpose models featured are: (1) DeepSeek-LLM-Chat 67B \citep{deepseek-llm}, (2) Qwen 72B \citep{qwen}, (3) SeaLLM-v2 7B \citep{seallm}, and (4) ChatGLM3 6B \citep{chatglm3}. In addition, several open-source models focus on superior mathematical capabilities, including: (5) InternLM2-Math 20B~\footnote{\url{https://github.com/InternLM/InternLM-Math}}, an InternLM2 derivative fine-tuned first with mathematics data and then using instruction tuning; (6) Math-Shepherd-Mistral 7B, which leverages PPO training \citep{schulman2017proximal} on top of Mistral 7B \citep{mistral} with a reward model based on process supervision; (7) the WizardMath suite \citep{wizardmath}, which strengthens mathematical reasoning for Mistral 7B and Llama-2 70B \citep{llama2} via the evolve-instruct methodology (a form of instruction tuning with AI-generated instructions) combined with PPO, predominantly using GSM8K and MATH as training sources; (8) MetaMath 70B \citep{metamath}, a version of Llama-2 70B fine-tuned on expanded versions of GSM8K and MATH datasets; (9) ToRA 34B \cite{tora}, which fine-tunes CodeLlama 34B for mathematical reasoning supported by tool integration; and (10) MAmmoTH 70B \citep{MathInstruct}, where Llama-2 70B is further refined on MathInstruct through instruction tuning.

Table \ref{tab:sft_rl_math} presents results for a setting in which models are restricted from using external tools. In this scenario, DeepSeekMath-Instruct 7B achieves robust step-by-step reasoning capabilities. Remarkably, on the challenging MATH benchmark, it exceeds the performance of all listed open-source competitors as well as most closed models (such as Inflection-2 and Gemini Pro) by at least 9\% absolute—this advantage is maintained over both larger models (e.g., Qwen 72B) and those specifically improved with math-focused reinforcement learning strategies (like WizardMath-v1.1 7B). DeepSeekMath-Instruct also matches the Chinese proprietary models GLM-4 and Baichuan-3 on the MATH dataset, but remains behind GPT-4 and Gemini Ultra in performance.

When models are evaluated with the option to leverage both natural language inference and programmatic tool integration, DeepSeekMath-Instruct 7B achieves nearly 60\% accuracy on the MATH dataset, establishing a new peak among open-source contenders. Across additional benchmarks, its results are on par with DeepSeek-LLM-Chat 67B—the prior leading system in this category despite being an order of magnitude larger.

\section{Reinforcement Learning}

\subsection{Group Relative Policy Optimization} Reinforcement learning (RL) has shown significant promise for further strengthening the mathematical reasoning abilities of LLMs post-Supervised Fine-Tuning (SFT) \citep{wang2023math,wizardmath}. Here, we propose Group Relative Policy Optimization (GRPO), our novel RL approach designed for efficiency and effectiveness.

\subsubsection{From PPO to GRPO} Proximal Policy Optimization (PPO) \citep{schulman2017proximal} is a widely adopted actor-critic reinforcement learning method during LLM fine-tuning phases \citep{ouyang2022training}. PPO seeks to maximize the following surrogate objective in tuning LLM parameters:

\begin{equation}
\footnotesize
    \mathcal{J}_{PPO}(\theta) = \mathbb{E}{[q \sim P(Q), o \sim \pi_{\theta_{old}}(O|q)]} \frac{1}{|o|} \sum_{t=1}^{|o|} \min \left[ \frac{\pi_\theta(o_{t} | q, o_{<t})}{\pi_{\theta_{old}}(o_{t} | q, o_{<t})} A_{t}, \text{clip} \left( \frac{\pi_\theta(o_{t} | q, o_{<t})}{\pi_{\theta_{old}}(

Here, $\pi_{\theta}$ denotes the current policy, while $\pi_{\theta_{old}}$ refers to the previous iteration of the policy model. The variables $q$ and $o$ represent questions and corresponding sampled outputs, drawn from the dataset and the old policy $\pi_{\theta_{old}}$, respectively. The constant $\varepsilon$ serves as a clipping hyper-parameter, introduced in PPO to ensure training stability. $A_t$ indicates the advantage at time $t$, which is typically computed using Generalized Advantage Estimation (GAE) \citep{gae} based on rewards $\{r_{\ge t}\}$ and an estimated value function $V_{\psi}$. Consequently, PPO necessitates training an auxiliary value function alongside the main policy model. To address potential overfitting to the reward model, a conventional technique is to include a per-token KL divergence penalty relative to a reference model in the reward computation for each token, as outlined in \citep{ouyang2022training}, namely,

\begin{equation}
    r_{t} = r_\varphi(q, o_{\le t}) - \beta \log\frac{\pi_{\theta}(o_{t}|q, o_{<t})}{\pi_{ref}(o_{t}|q, o_{<t})},
\label{eq:PPO-reward}
\end{equation}

where $r_\varphi$ represents the reward model, $\pi_{ref}$ is the reference policy—often initialized as the SFT model—and $\beta$ determines the magnitude of the KL penalty.

The incorporation of a value function in PPO usually demands a model of a size similar to the main policy, resulting in significant resource and computational overhead. Furthermore, during RL fine-tuning, the value function acts as a baseline to reduce the variance of advantage estimation. In large language model (LLM) contexts, it is common practice for the reward model to only evaluate the final token of an output sequence, which may pose challenges for accurately training a value function at the token level. To circumvent this, we introduce Group Relative Policy Optimization (GRPO), as illustrated in Figure \ref{fig:grpo}. GRPO eliminates the necessity for separate value function estimation as used in PPO; instead, it adopts the mean reward from multiple outputs generated in response to the same question as the baseline. In practice, for each prompt $q$, a collection of outputs $\{o_1, o_2, \cdots, o_G\}$ is sampled from the old policy $\pi_{\theta_{old}}$, and the optimization of the policy model is carried out by maximizing the following objective:

\begin{equation}
\footnotesize
\begin{split}
    \mathcal{J}_{GRPO}(\theta) &= \mathbb{E}{[q \sim P(Q), \{o_i\}_{i=1}^G \sim \pi_{\theta_{old}}(O|q)]}  \\
    & \frac{1}{G}\sum_{i=1}^G\frac{1}{|o_i|} \sum_{t=1}^{|o_i|} \left\{ \min \left[ \frac{\pi_\theta(o_{i,t} | q, o_{i,<t})}{\pi_{\theta_{old}}(o_{i,t} | q, o_{i,<t})} \hat{A}_{i,t}, \text{clip} \left( \frac{\pi_\theta(o_{i,t} | q, o_{i,<t})}{\pi_{\theta_{old}}(o_{i,t} | q, o_{i,<t})}, 1 - \varepsilon, 1 + \varepsilon \right)  \hat{A}_{i,t} \right] - \beta \mathbb{D}_{KL}\left[\pi_{\theta} || \pi_{ref}\right]\right\} ,
\end{split}
\label{eq:GRPO-obj}
\end{equation}

with $\varepsilon$ and $\beta$ as hyper-parameters. The advantage $\hat{A}_{i,t}$ is estimated utilizing only the rewards within the respective group of sampled outputs; further explanation is provided in subsequent sections. GRPO’s use of group-relative advantage computation naturally aligns with the comparative design of reward models, as such models are commonly trained on data comprising comparisons among responses to the same question. Notably, rather than incorporating the KL penalty within the reward, GRPO enforces regularization by appending the KL divergence between the policy under training and the reference policy directly to the loss function. This simplifies the calculation of $\hat{A}_{i,t}$. Unlike the KL regularization term in (\ref{eq:PPO-re

\begin{algorithm}[t]
  \small
  \caption{Iterative Group Relative Policy Optimization}
  \textbf{Input} initial policy model $\pi_{\theta_{\text{init}}}$; reward models $r_\varphi$; task prompts $\mathcal{D}$; 
  hyperparameters $\varepsilon$, $\beta$, $\mu$
  \begin{algorithmic}[1]
    \State Set $\pi_\theta \leftarrow \pi_{\theta_{\text{init}}}$
    \For{iteration = 1 to $I$}
      \State Set reference policy $\pi_{ref} \leftarrow \pi_{\theta}$
      \For{step = 1 to $M$}
        \State Select a minibatch $\mathcal{D}_b$ from $\mathcal{D}$
        \State Assign $\pi_{\theta_{old}} \leftarrow \pi_{\theta}$
        \State For every question $q$ in $\mathcal{D}_b$, generate $G$ outputs $\{o_i\}_{i=1}^G \sim \pi_{\theta_{old}} (\cdot \mid q) $
        \State For each generated output $o_i$, evaluate reward values $\{r_i\}_{i=1}^{G}$ by applying $r_{\varphi}$
        \State Estimate group-relative advantages $\hat{A}_{i,t}$ at each token $t$ in output $o_i$
        \For{GRPO step = 1 to $\mu$}
          \State Update policy parameters $\pi_\theta$ by optimizing the GRPO objective (see Equation~\ref{eq:GC-GRPO})
        \EndFor
      \EndFor
      \State Periodically refine $r_\varphi$ by training with a replay buffer.
    \EndFor
  \end{algorithmic}
  \textbf{Output} $\pi_\theta$
  \label{alg:iter-grpo}
\end{algorithm}

\subsubsection{Outcome Supervision RL with GRPO} In the formal setting, for each input question $q$, a collection of $G$ responses $\{o_1, o_2, \cdots, o_G\}$ are sampled from $\pi_{\theta_{old}}$. The reward model scores each output, producing a corresponding vector of $G$ reward values $\mathbf{r}=\{r_1, r_2, \cdots, r_G\}$. These rewards are standardized by subtracting the mean and dividing by the standard deviation within the group. Under outcome supervision, the final normalized reward is provided for each complete output $o_i$ and is assigned as the advantage $\hat{A}_{i, t}$ for every token in $o_i$: that is, $\hat{A}_{i, t} = \widetilde{r}_i = \frac{r_i- {\rm mean}(\mathbf{r})}{{\rm std}(\mathbf{r})}$. The policy is then trained to maximize the group-relative objective specified in Equation~(\ref{eq:GRPO-obj}).

\subsubsection{Process Supervision RL with GRPO} Because outcome supervision only gives terminal rewards, it can lack sufficient granularity for guiding the model in intricate mathematical scenarios. Adopting the approach of \cite{wang2023math}, we also implement process supervision, which delivers feedback at every reasoning step. Specifically, given question $q$ and a batch of $G$ sampled outputs $\{o_1, o_2, \cdots, o_G\}$, a process reward model evaluates each reasoning stage within the outputs, generating reward sets $\mathbf{R} = \{ \{r_1^{index(1)},\cdots,r_1^{index(K_1)}\}, \cdots,  \{r_G^{index(1)},\cdots,r_G^{index(K_G)}\} \}$, where $index(j)$ is the index of the terminal token for the $j$th step, and $K_i$ is the step count for $o_i$. The set of rewards is normalized by centering and scaling: $\widetilde{r}_i^{index(j)} = \frac{r_i^{index(j)} - {\rm mean(\mathbf{R})}}{{\rm std(\mathbf{R})}}$. Advantage for each token is then given as the sum of future normalized rewards, i.e., $\hat{A}_{i, t} = \sum_{index(j) \ge t} \widetilde{r}_i^{index(j)}$. Training then proceeds by maximizing the same objective as in Equation~(\ref{eq:GRPO-obj}).

\subsubsection{Iterative RL with GRPO} As reinforcement learning advances, the reward model used initially may become inadequate for guiding the continually improving policy model. To address this, we implement iterative RL with GRPO. In this framework, illustrated in Algorithm \ref{alg:iter-grpo}, we iteratively generate updated training data for the reward model using samples produced by the current policy model. The reward model is then continually fine-tuned with a replay mechanism that incorporates $10\%$ of past data. Afterward, the reference model is replaced with the latest policy model, which is then further trained using the most recent version of the reward model.

\subsection{Training and Evaluating DeepSeekMath-RL}

Our RL experiments use DeepSeekMath-Instruct 7B as the foundation. The RL training set consists of chain-of-thought style questions related to GSM8K and MATH, sourced from the SFT data, totaling approximately $144$K questions. Other SFT items are excluded, enabling analysis of RL effects on evaluation benchmarks that are not represented during RL training. Construction of the reward model’s training set follows the methodology in \citep{wang2023math}. The initial reward model is trained on DeepSeekMath-Base 7B, with a learning rate of $2\mathrm{e}{-5}$. The GRPO training utilizes a policy model learning rate of $1\mathrm{e}{-6}$ and a KL coefficient of $0.04$. For each prompt, $64$ completions are sampled. The sequence length is capped at $1024$, and the training batch size is also $1024$. The policy model undergoes one update step after each exploration period. Evaluation is performed on the same benchmarks as DeepSeekMath-Instruct 7B. For DeepSeekMath-RL 7B, tasks from GSM8K and MATH with chain-of-thought formats are considered in-domain, whereas all other benchmarks are classified as out-of-domain.

Table \ref{tab:sft_rl_math} presents a comparison of open- and closed-source models in terms of their capabilities in both chain-of-thought and tool-integrated reasoning on benchmarks in English and Chinese. Our observations reveal that: 1) DeepSeekMath-RL 7B achieves accuracies of 88.2\% on GSM8K and 51.7\% on MATH under chain-of-thought prompting. This positions it ahead of all open-source models within the 7B to 70B parameter range, as well as surpassing most closed-source alternatives. 2) Notably, DeepSeekMath-RL 7B is exclusively trained on chain-of-thought instruction tuning datasets for GSM8K and MATH, initializing from DeepSeekMath-Instruct 7B. In spite of this limited training data, DeepSeekMath-RL 7B consistently exceeds the performance of DeepSeekMath-Instruct 7B across all evaluation metrics, underscoring the advantage conferred by reinforcement learning.

\section{Discussion}
In this part, we detail insights gleaned from our investigations in pre-training and reinforcement learning.

\subsection{Lessons Learnt in Pre-Training}

We begin by outlining our experience and observations related to the pre-training stage. Unless noted otherwise, the configurations align with those described in Section \ref{sec:quality-policy}. Additionally, the term DeepSeekMath Corpus in this context denotes an 89B-token dataset, which is the product of our data curation’s second cycle.

\subsubsection{Code Training Benefits Mathematical Reasoning}

The belief that code training enhances models’ reasoning ability is widespread, though often lacking empirical validation. We seek to contribute concrete evidence within the mathematical domain, showing that code-based training augments models’ mathematical reasoning, regardless of whether tool integration is present.

To evaluate the impact of code training on mathematical reasoning proficiency, we adopted both two-stage and one-stage training paradigms in our experimentation:

\textbf{Two-Stage Training}

\begin{itemize}[topsep=0pt]
    \item \textbf{Code Training for 400B Tokens $\rightarrow$ Math Training for 150B Tokens}: The DeepSeek-LLM 1.3B model is initially exposed to 400B code tokens, which is subsequently followed by fine-tuning on 150B math tokens.
    \item \textbf{General Training for 400B Tokens $\rightarrow$ Math Training for 150B Tokens}: For comparison, we also examine a training regimen where the first phase consists of 400B general tokens (sampled from a comprehensive general corpus developed by DeepSeek-AI) instead of code tokens. This approach is aimed at evaluating the comparative impact of code versus general token pre-training on the improvement of mathematical reasoning skills.
\end{itemize}

\textbf{One-Stage Training}

\begin{itemize}[topsep=0pt]
    \item \textbf{Math Training for 150B Tokens}: Here, DeepSeek-LLM 1.3B is trained solely on 150B math tokens from the outset.
    \item \textbf{Training on a mixture of 400B Code Tokens and 150B Math Tokens}: Recognizing that math training following code training can impair performance on coding tasks, we further explore whether jointly training on both code and math tokens in a single stage can both enhance mathematical reasoning and counteract the effects of catastrophic forgetting.
\end{itemize}

\paragraph{Results}
Tables \ref{tab:code-math} and \ref{tab:code-math-general-eval} present the model's performance across different training paradigms.

Pre-training on code has a pronounced positive effect on program-aided mathematical reasoning, whether utilizing sequential (two-stage) or combined (one-stage) strategies. According to Table \ref{tab:code-math}, employing a two-stage procedure where code token training precedes math training greatly enhances performance on GSM8K and MATH datasets that require Python-based solutions. Additional gains are observed with math-specific fine-tuning in the subsequent phase. For the one-stage setup, simultaneous exposure to both code and math tokens notably reduces the catastrophic forgetting that can follow separate staged training, and also results in improvements to both coding proficiency (see Table \ref{tab:code-math-general-eval}) and math reasoning with program assistance (refer to Table \ref{tab:code-math}).

Furthermore, code token pre-training strengthens mathematical reasoning even in contexts lacking external tool use. Within the two-stage regime, code token training alone already brings about noticeable advances. It further accelerates and strengthens later math-specific training, ultimately delivering the highest observed performance. In contrast, single-stage mixed training with both code and math tokens detracts from mathematical reasoning performance in settings where tool use is not permitted. This might be attributable to DeepSeek-LLM 1.3B’s constrained model size, which may be insufficient to concurrently absorb large-scale code and mathematical datasets.

\subsubsection{ArXiv Papers Appear Unhelpful for Advancing Mathematical Reasoning}

Although ArXiv papers are a standard constituent of math pre-training datasets \citep{minerva,gpt-f,llemma,mathpile}, their precise contribution to enhancing mathematical reasoning abilities has yet to be thoroughly investigated. Somewhat surprisingly, our empirical findings indicate that ArXiv papers do not effectively support improvements in mathematical reasoning. We performed systematic experiments across different model sizes—including DeepSeek-LLM 1.3B and DeepSeek-Coder-Base-v1.5 7B \citep{deepseek-coder}—by leveraging ArXiv-based corpora subjected to distinct processing approaches:

\begin{itemize}[topsep=0pt]
    \item \textbf{MathPile} \citep{mathpile}: an 8.9B-token dataset crafted using various heuristics for data cleaning and filtering, where over 85\% of the content originates from scientific ArXiv papers;
    \item \textbf{ArXiv-RedPajama} \citep{redpajama}: comprising all available ArXiv LaTeX files with supplementary material such as preambles, comments, macros, and references excluded, totaling 28.0B tokens.
\end{itemize}

To assess the impact, we conducted isolated pre-training for DeepSeek-LLM 1.3B on 150B tokens and for DeepSeek-Coder-Base-v1.5 7B on 40B tokens for each ArXiv corpus. Results show that models trained exclusively on ArXiv content fail to deliver tangible gains, and may even experience performance drops, on a range of mathematical reasoning benchmarks used in this study. Evaluations include quantitative reasoning datasets such as GSM8K and MATH (refer to Table \ref{tab:arxiv-cot}), multiple-choice assessments like MMLU-STEM (see Table \ref{tab:arxiv-cot}), as well as formal mathematics benchmarks including miniF2F (see Table \ref{tab:arxiv_atp}).

Nonetheless, these findings are not without caveats. Important aspects remain unaddressed in our analysis, including:

\begin{itemize}[topsep=0pt]
    \item The potential influence of ArXiv data on math-adjacent tasks absent from our study, such as the transformation of formal theorems and proofs into their informal counterparts;
    \item The outcomes when ArXiv material is integrated with other data sources;
    \item Whether positive effects of ArXiv pre-training emerge at scales beyond those investigated.
\end{itemize}

Consequently, additional research is warranted, which we propose as avenues for future work.

\subsection{Insights of Reinforcement Learning}
\subsubsection{Towards to a Unified Paradigm}
Here, we outline a unified framework to systematically study various training techniques—such as SFT, RFT, DPO, PPO, GRPO—and present experiments to dissect key elements within this framework. In general, the gradient of a particular training method with respect to its parameter $\theta$ can be represented as:

\begin{equation}
    \nabla_{\theta}\mathcal{J}_{\textcolor{red}{\mathcal{A}}}(\theta) = \mathbb{E}[\underbrace{(q,o) \sim \textcolor{red}{\mathcal{D}}}_{Data \ Source}]\left( \frac{1}{|o|} \sum_{t=1}^{|o|}  \underbrace{GC_{{\mathcal{A}}}(q, o, t, \textcolor{red}{\pi_{{rf}}})}_{Gradient \ Coefficient}  \nabla_{\theta}\log \pi_{\theta}(o_t | q, o_{<t})\right).
\label{eq:objective}
\end{equation}

This framework comprises three essential elements: (1) the \textit{Data Source $\mathcal{D}$}, specifying the origin of training examples; (2) the \textit{Reward Function $\pi_{{rf}}$}, responsible for providing reward signals during training; and (3) the \textit{Algorithm $\mathcal{A}$}, which translates both the training data and reward signals into the gradient coefficient $GC$, governing the extent of reward or penalty applied. Using this general formulation, we examine several notable approaches:

\begin{itemize}
    \item  \textbf{Supervised Fine-tuning (SFT)}: This approach adapts a pretrained model by additional training on SFT datasets curated by humans.
    \item \textbf{Rejection Sampling Fine-tuning (RFT)}: RFT further tunes the model resulting from SFT by training on filtered model responses to SFT prompts, where only outputs deemed correct are retained.
    \item \textbf{Direct Preference Optimization (DPO)}: DPO continues to improve the SFT model by employing pairwise preference-based DPO loss on an expanded set of outputs drawn from the SFT model.
    \item \textbf{Online Rejection Sampling Fine-tuning (Online RFT)}: Unlike RFT, Online RFT uses the SFT model for initialization and iteratively fine-tunes it with new outputs generated in real time by the current policy model.
    \item \textbf{PPO/GRPO}: These algorithms start with the SFT model as initialization and apply reinforcement via outputs dynamically generated by the continually updated policy model.
\end{itemize}

A summary of these methods' components is presented in Table \ref{tab:data_policy}. For a more comprehensive discussion of the derivation, please see Appendix \ref{app:analysis-rl}.

\paragraph{Observation about Data Source} We categorize data sources into online and offline sampling. Online sampling refers to data collected from the real-time exploratory actions of the current training policy model, whereas offline sampling refers to data obtained from the sampling outcomes of the initial SFT model. Notably, RFT and DPO employ an offline approach, whereas Online RFT and GRPO utilize an online strategy.

Figure \ref{fig:rl-analysis} illustrates that Online RFT achieves substantial gains over RFT on both evaluated benchmarks. Initially, the performance of Online RFT is nearly identical to that of RFT; however, as training proceeds, Online RFT establishes a clear and consistent lead, underscoring the effectiveness of online data collection. This observation is expected: in early training phases, the actor and the SFT model remain closely aligned, so their sampled outputs differ only marginally. As training progresses, the divergence in samples from the actor increases, and continually sampling fresh data in real time confers additional benefits.

\paragraph{Observation about Gradient Coefficient} To update model parameters, the algorithm transforms the reward signal into a gradient coefficient. In our experiments, we distinguish between two reward types: `Rule', which evaluates the quality of a response according to its factual correctness, and `Model', where a separately trained reward model assigns a score to each response (with its training data based on rule-based evaluations). Equations \ref{eq:GC-RFT} and \ref{eq:GC-GRPO} reveal an important distinction between GRPO and Online RFT. GRPO dynamically adjusts the gradient coefficient for each response according to its reward, enabling varying degrees of reinforcement or penalization. Conversely, Online RFT does not make this differentiation; it treats all correct responses equally and does not explicitly penalize incorrect responses.

As illustrated in Figure \ref{fig:rl-analysis}, GRPO outperforms online RFT, underscoring the effectiveness of modulating positive and negative gradient coefficients. Additionally, GRPO+PS achieves better results than GRPO+OS, suggesting that leveraging step-sensitive, granular gradient coefficients provides further advantages. We also investigate iterative reinforcement learning by performing two iterations in our experiments. Figure \ref{fig:iter-rl} shows that performance is notably enhanced through iterative RL, with the most substantial gains occurring in the initial iteration.

\subsubsection{Why RL Works?} This study implements reinforcement learning on a selected subset of instruction tuning data, which yields marked improvements over models trained solely with instruction tuning. To further clarify the reasons behind RL’s effectiveness, we assess both Pass@K and Maj@K accuracies for the Instruct and RL models on two distinct benchmarks. As depicted in Figure \ref{fig:maj-pass}, RL contributes to a noticeable uplift in Maj@K while leaving Pass@K largely unaffected. This outcome suggests that RL’s strength lies in increasing the likelihood of correct answers within the TopK outputs, rather than enhancing the model’s core problem-solving abilities. Similarly, \citep{wang2023making} observed a \textbf{misalignment problem} when reasoning with SFT models, revealing that reasoning accuracy can be improved by employing various preference alignment strategies \citep{yuan2023rrhf,song2023preference,wang2023making}.

\subsubsection{How to Achieve More Effective RL?}
\label{sec:effec_rl}
Our findings show that reinforcement learning is highly effective for mathematical reasoning problems. Moreover, we propose a unifying framework for understanding prominent training techniques, wherein each approach is interpreted as a direct or approximate form of RL. According to Equation \ref{eq:objective}, this framework consists of three critical components: Data Source, Algorithm, and Reward Function. We outline several prospective research directions for each.

\paragraph{Data Source} The data source is fundamental to every training approach. In RL, this refers to the set of unlabeled questions paired with outputs sampled from the policy model. In our work, we rely exclusively on questions from the instruction tuning stage and utilize a basic nucleus sampling strategy for output generation. This methodological choice may explain why our RL pipeline chiefly improves Maj@K. For future studies, we intend to expand our RL pipeline to include out-of-distribution prompts and to integrate \textbf{advanced sampling (decoding) strategies}, such as tree-search-based methods \citep{yao2023tree}. In addition, employing \textbf{efficient inference techniques} \citep{xia-etal-2023-speculative,leviathan2023fast,kwon2023efficient,xia2024unlocking}—which critically affect the policy model’s exploration efficiency—remains an important direction to pursue.

\paragraph{Algorithms} Algorithms utilize the available data and the reward signal to compute the gradient coefficient, thereby updating the model parameters. As implied by Equation \ref{eq:objective}, contemporary approaches largely \textbf{TRUST} the guidance from the reward function, adjusting the conditional probability of certain tokens upward or downward according to the received signal. Nonetheless, there is no guarantee that the reward signal maintains high fidelity in all cases, particularly when dealing with highly intricate tasks. For instance, even meticulously curated datasets like PRM800K \citep{lightman2023let}, which are annotated by skilled professionals, exhibit roughly 20\% inaccurate labels\footnote{\url{https://github.com/openai/prm800k/issues/12\#issuecomment-1728491852}}. This highlights the necessity of developing reinforcement learning algorithms that can maintain robustness in the presence of noisy reward signals. We posit that the adoption of \textbf{WEAK-TO-STRONG} alignment strategies \citep{burns2023weak} has the potential to transform the core paradigms of learning algorithms.

\paragraph{Reward Function} The reward function acts as the principal provider of training signals. Within the context of RL, it is commonly realized as a neural reward model. We identify three critical research avenues for reward models: 1) \textbf{Improving the generalization capacity of the reward model.} It is imperative for the reward model to generalize effectively to unforeseen queries and novel decoding results; failure to do so may result in RL only consolidating the output distribution of LLMs instead of enhancing their underlying abilities; 2) \textbf{Capturing and quantifying reward model uncertainty.} Incorporating uncertainty estimation could offer a means of bridging the gap between underperforming reward models and weak-to-strong alignment methodologies; 3) \textbf{Designing scalable, high-fidelity process reward models} capable of supplying nuanced, detailed signals throughout the reasoning process \citep{lightman2023let,wang2023math}.

\section{Conclusion, Limitation, and Future Work}
This work introduces DeepSeekMath, which surpasses all open-source counterparts on the competition-level MATH benchmark and attains results comparable to those of proprietary models. DeepSeekMath~leverages DeepSeek-Coder-v1.5 7B as its foundation and is subjected to an extended training regimen comprising 500B tokens, notably including 120B math-related tokens extracted from Common Crawl. Through comprehensive ablation studies, we demonstrate that web-sourced data constitutes a valuable source of high-quality mathematical content, whereas contributions from arXiv are less substantial than anticipated. We propose Group Relative Policy Optimization (GRPO), an adaptation of Proximal Policy Optimization (PPO), which significantly enhances mathematical reasoning performance while reducing memory requirements. Experimental findings reveal that GRPO delivers consistent improvements even as DeepSeekMath-Instruct 7B achieves top-tier results on benchmarks. Additionally, we present a unified analytical framework for various techniques and identify multiple promising directions for advancing reinforcement learning efficacy.

Despite DeepSeekMath's strong results on quantitative reasoning tasks, its proficiency in geometry and theorem-proving lags behind that of proprietary models. Notably, in our preliminary evaluations, the model was unable to effectively address problems involving triangles and ellipses, suggesting that biases in pre-training and fine-tuning data selection may be present. Furthermore, due to its model size, DeepSeekMath~does not match GPT-4 in few-shot learning scenarios; while GPT-4's performance benefits from few-shot inputs, DeepSeekMath~exhibits similar outcomes across zero-shot and few-shot tests. For future work, we aim to further enhance our data selection processes to assemble an even higher-quality pre-training dataset. Moreover, we will investigate additional strategies (as discussed in Section \ref{sec:effec_rl}) for optimizing reinforcement learning approaches in LLM development.